\section*{الفصل \ref{ch:b2:vertebrate-development} --- النمو الجنيني عند الفقاريات}

\begin{solution}{exo:b2:vertebrate-development:1}
\omterm{def:b2:vertebrate-development:egg}{التفلج}: تقسمات خيطية سريعة بلا نمو، تنتج \omterm{def:b2:vertebrate-development:blastula}{الأريمة}.
\omterm{def:b2:vertebrate-development:blastula}{والأريمة}: كرة جوفاء أو قرص من خلايا صغيرة حول
\omterm{def:b2:vertebrate-development:blastula}{جوف الأريمة}. والتمعّي: حركات الخلايا التي تُدخل الأديمين المتوسط
والباطن، فتنتج المعيولة الثلاثية الطبقات بجوف معوي
\omterm{prop:b2:vertebrate-development:gastrulation}{وحبل ظهري}. والتعصّب: انطواء الأديم الظاهر الظهري
أنبوبًا عصبيًا، مع تقطع القُطَيعات،
فينتج العصبولة بالمخطط التنظيمي الفقاري.
\end{solution}

\begin{solution}{exo:b2:vertebrate-development:2}
الأديم الظاهر: البشرة، والدماغ والنخاع الشوكي، والأعصاب المحيطية
(من \omterm{prop:b2:vertebrate-development:neurulation}{العرف العصبي})، والعدسة وخلايا الصباغ. والأديم المتوسط: \omterm{prop:b2:vertebrate-development:gastrulation}{الحبل
الظهري}، والعضل، والعظم والغضروف، والكلية، والقلب والدم، والأدمة.
والأديم الباطن: بطانة المعي، والكبد، والمعثكلة، والرئتان، والدرقية.
\end{solution}

\begin{solution}{exo:b2:vertebrate-development:3}
\omterm{def:b2:vertebrate-development:egg}{مح} قليل: \omterm{def:b2:vertebrate-development:egg}{تفلج} تام شبه متساوٍ وتمعٍّ
بالانقلاب عبر \omterm{prop:b2:vertebrate-development:gastrulation}{ثغرة أريمية} مستديرة في كرة جوفاء (والضفدع:
تام لكن غير متساوٍ، وخلاياه النباتية المحية كبيرة بطيئة). \omterm{def:b2:vertebrate-development:egg}{ومح}
كثير (الكتكوت): \omterm{def:b2:vertebrate-development:egg}{تفلج} محصور في قرص فوق \omterm{def:b2:vertebrate-development:egg}{المح}، و
تمعٍّ عبر شق هو الخط البدئي، في القرص
المسطح، والطبقات تنتشر فوق \omterm{def:b2:vertebrate-development:egg}{المح} بدل أن تحيط
بجوف.
\end{solution}

\begin{solution}{exo:b2:vertebrate-development:4}
الاستحثاث تغيير مصير نسيج أهل بإشارة
من نسيج مجاور. فالخلايا النباتية تستحث أديمًا متوسطًا في
المنطقة الهامشية فوقها؛ \omterm{def:b2:vertebrate-development:induction}{والمنظّم} (أديم شفة الثغرة الظهرية المتوسط) يستحث
الصفيحة العصبية في الأديم الظاهر الظهري؛ والحويصلة البصرية تستحث
العدسة في أديم الرأس الظاهر.
\end{solution}

\begin{solution}{exo:b2:vertebrate-development:5}
$2^{12} = 4096$ خلية نصف قطرها $1/2^{4} = \qty{62.5}{\micro m}$؛
والسطح مضروب في $2^{4} = 16$. والسطح الزائد هو
غشاء الظهارات التي سيطويها التمعّي جلدًا
ومعيًا، والمساحة التي تتبادل عبرها \omterm{def:b2:vertebrate-development:blastula}{الأريمة} مع
محيطها.
\end{solution}

\begin{solution}{exo:b2:vertebrate-development:6}
$10^{-4}\times 2^{k} = 0.4$: ومنه $2^{k} = 4000$، أي $k = 12$. ومع أربعة
أمثال DNA تبدأ النسبة عند $4\times 10^{-4}$ وتبلغ $0.4$
عند $2^{k} = 1000$: أي الانقسام العاشر، أي أبكر باثنين.
\end{solution}

\begin{solution}{exo:b2:vertebrate-development:7}
$k = \ln 2/600 = \qty{1.16e-3}{s^{-1}}$؛ و$\lambda = \sqrt{2\times
10^{-12}/1.16\times 10^{-3}} = \qty{42}{\micro m}$. والعتبتان عند
$x = \lambda\ln 2 = \qty{29}{\micro m}$ و$\lambda\ln 10 =
\qty{96}{\micro m}$: أي نطق من 29 و67 وما بقي.
\end{solution}

\begin{solution}{exo:b2:vertebrate-development:8}
مضاعفة المنبع تزيح الحدين إلى الخارج بمقدار $\lambda\ln
2 = \qty{29}{\micro m}$: فيتضاعف النطاق الأول، ويحفظ الثاني
عرضه. ومضاعفة $k$ تقسم $\lambda$ على $\sqrt 2$
(أي \qty{29}{\micro m})، ولمّا كان $c_0 = j/\sqrt{Dk}$، فهي تقسم $c_0$ على
$\sqrt 2$ أيضًا: فينتقل الحدان إلى $\lambda'\ln(c_0'/T) =
29\times(0.69 - 0.35) = \qty{10}{\micro m}$ و$29\times(2.30 -
0.35) = \qty{57}{\micro m}$ --- فينكمش النطاقان معًا.
\end{solution}

\begin{solution}{exo:b2:vertebrate-development:9}
طُعّمت شفة ظهرية من معيولة سمندل غير مصبوغ على
الجانب البطني لمعيولة مصبوغة؛ فتدحرجت إلى الداخل وكوّن المضيف
جنينًا ثانيًا --- أنبوبًا عصبيًا وقُطَيعات ومعيًا --- موصولًا بالأول.
وبيّن فرق التصبغ أن المحور الثاني كان من خلايا المضيف،
فالطُّعم إذن استحثها بدل أن يبني المحور
بنفسه. وقد استبعد ذلك التمايز الذاتي للطُّعم تفسيرًا،
وبيّن أن منطقة صغيرة تحمل الإشارة التي
تنظّم المحور كله.
\end{solution}

\begin{solution}{exo:b2:vertebrate-development:10}
(أ) جلد بطن: فالأديم الظاهر الباكر أهل لكنه لم يُستحث بعد،
فيتبع محيطه الجديد. (ب) نسيج عصبي: فقد استُحث
بحلول المعيولة المتأخرة وصار متحددًا. (ج) صفيحة عصبية:
فأديم البطن الظاهر في المعيولة المتأخرة لا يزال أهلًا، \omterm{prop:b2:vertebrate-development:gastrulation}{والحبل الظهري}
يستحثه. (د) لا بنى ظهرية: بل كرة من نسيج بطني، إذ
لا شيء يستحث الأديم المتوسط على مصائر ظهرية ولا الأديم الظاهر على
مصائر عصبية.
\end{solution}

\begin{solution}{exo:b2:vertebrate-development:11}
عند \qty{18}{\celsius}: $90 + 11\times 30 = \qty{420}{min}$، أي 7 ساعات.
وعند \qty{10}{\celsius}: \qty{17.5}{h}. ويقع الانتقال عند
الانقسام الثاني عشر في الحالتين لأن الذي يطلقه هو
نسبة DNA إلى السيتوبلازم، وهي تتوقف على عدد الانقسامات
لا على المدة التي استغرقتها: فالساعة تعدّ المضاعفات لا
الساعات.
\end{solution}

\begin{solution}{exo:b2:vertebrate-development:12}
في النمو التنظيمي تكتسب الخلايا مصائرها بالاستحثاث
من الجيران، فيستطيع الجنين أن يعيد البناء بعد فقد (توأمة من
بيضة مشقوقة، وتجدد منطقة مطعّمة)؛ أما في النمو
الفسيفسائي فتُورث المصائر مع محددات سيتوبلازمية
موضعية، فلا تصنع كل فلجة إلا جزءها ويصنع الجنين المشقوق
يرقتين نصفيتين. والفرق فرق في الدرجة: فهلال الضفدع
الرمادي محدد، والأجنة الفسيفسائية تستعمل
الاستحثاثات لاحقًا؛ فكل جنين يمزج الوراثة والمحادثة،
والفقاريات تتكئ على المحادثة زمنًا أطول.
\end{solution}

\begin{solution}{pb:b2:vertebrate-development:1}
\textbf{1.} البيضة $\tfrac43\pi(0.7)^{3} = \qty{1.44}{mm^3}$؛ والنواة
$\tfrac43\pi(0.03)^{3} = \qty{1.1e-4}{mm^3}$؛ والنسبة $8\times
10^{-5}$.
\textbf{2.} $2^{k} = 4000$: ومنه $k = 12$، أي 4096 خلية.
\textbf{3.} $75 + 11\times 30 = \qty{405}{min}$، أي نحو 6.75 ساعة.
\textbf{4.} نصف القطر $0.7/2^{4} = \qty{44}{\micro m}$؛ والنسبة $8\times
10^{-5}\times 4096 = 0.32$ --- وهي نسبة خلية عادية: فقد
لحقت النوى بالسيتوبلازم.
\textbf{5.} $2^{4} = 16$.
\textbf{6.} نحو $4095\,c$ من DNA؛ وتحوي البيضة ما يكفي من
الهستونات والبوليميرازات والنوكليوتيدات المخزونة لاثني عشر انقسامًا، فلا
حاجة إلى استنساخ.
\textbf{7.} تشتعل المورّثات اللاقحية، وتبطئ الانقسامات وتفقد
تزامنها، وتصير الخلايا متحركة. وحقن DNA زائد يقدّم
الانتقال بعدد الانقسامات الذي يفيض به DNA،
وهو ما لا يستطيع عدّاد انقسامات أو عدّاد زمن أن يفسره.
\textbf{8.} مقابل نقطة دخول النطفة، عند الهلال الرمادي
الذي ثبّته دوران قشرة البيضة بعد الإلقاح.
\textbf{9.} الأديم المتوسط الظهري أولًا (الصفيحة أمام الحبلية، ثم \omterm{prop:b2:vertebrate-development:gastrulation}{الحبل الظهري})
على سقف \omterm{prop:b2:vertebrate-development:gastrulation}{المعي البدئي} في الخط الناصف؛ ثم الأديم المتوسط الجانبي
(القُطَيعات، والصفيحة الجانبية) بجانبه؛ ويبطن الأديم الباطن
\omterm{prop:b2:vertebrate-development:gastrulation}{المعي البدئي}؛ وينتشر الأديم الظاهر على الخارج بالانتشار السطحي.
\textbf{10.} $\qty{2000}{\micro m}/\qty{600}{min} = \qty{3.3}{\micro
m/min}$: أي ثلاثة أمثال أرومة ليفية --- فالصفيحة المتناسقة تتحرك
أسرع من خلية منفردة.
\textbf{11.} تحتاج حركات الخلايا إلى نواتج مورّثات لاقحية (جزيئات
التصاق جديدة، وقدرة على الحركة) وإلى دورات أبطأ؛ وقبل
الانتقال لا تستطيع الخلايا إلا الانقسام.
\textbf{12.} لكل حبليّ \omterm{prop:b2:vertebrate-development:gastrulation}{حبل ظهري} في طور ما، و
تبني الفقاريات العمود الفقري حوله؛ أما \omterm{prop:b2:vertebrate-development:neurulation}{الأنبوب العصبي}
فيستحثه \omterm{prop:b2:vertebrate-development:gastrulation}{الحبل الظهري} وهو نتيجة المخطط لا أصله.
\textbf{13.} دون تقلص الخلايا القارورية لا تتكون \omterm{prop:b2:vertebrate-development:gastrulation}{ثغرة أريمية}،
ولا يتدحرج شيء إلى الداخل، ويبقى الجنين كرة بطبقاته في
مواضعها: فلا معي ولا \omterm{prop:b2:vertebrate-development:gastrulation}{حبل ظهري} ولا محور.
\textbf{14.} $k = \ln 2/1200 = \qty{5.8e-4}{s^{-1}}$؛ و$\lambda =
\sqrt{10^{-12}/5.8\times 10^{-4}} = \qty{42}{\micro m}$.
\textbf{15.} $x_1 = 42\ln 2.5 = \qty{38}{\micro m}$؛ و$x_2 = 42\ln 12.5
= \qty{105}{\micro m}$.
\textbf{16.} نطق من \qty{38}{\micro m} (أي 2.5 خلية)، و\qty{67}{\micro
m} (أي 4.5 خلية)، وما بقي.
\textbf{17.} $x_1 = 42\ln 1.25 = \qty{9}{\micro m}$؛ و$x_2 = 42\ln 6.25
= \qty{77}{\micro m}$: فينتقل الحدان إلى الداخل بمقدار $\lambda\ln 2 =
\qty{29}{\micro m}$؛ وينكمش النطاق الأول من 38 إلى 9، ويحفظ النطاق
الأوسط عرضه، وينمو النطاق الخارجي.
\textbf{18.} المنقولة قبل القراءة تصير المصير أ --- فهي تقرأ
التركيز حيث هي الآن؛ والمنقولة بعدها تبقي المصير ب --- فقد
تحددت.
\textbf{19.} $x_T = \lambda\ln(c_0/T)$: فتغير المنبع بعامل
$f$ يحرك كل حد بمقدار $\lambda\ln f$، فخطأ بعامل اثنين
يزيح النميط بمقدار $0.7\lambda$ فقط، وهو كسر من حقل
يمتد عدة $\lambda$ --- فالنميط لا يكاد يتأثر
بمقدار المورفوجين بالضبط.
\textbf{20.} شفة ظهرية إلى البطن: محور ثانٍ. ومنطقة هامشية
بطنية إلى الظهر: تُعاد تحديدها بمحيطها فتصنع
نسيجًا ظهريًا؛ ولا شيء زائد.
\textbf{21.} في العصبولة لم يعد الأديم الظاهر أهلًا: فيصنع
الطُّعم حبلًا ظهريًا ولا يستحث أنبوبًا عصبيًا ثانيًا.
\textbf{22.} الشق على امتداد مستوي الهلال: جنينان كاملان؛
والشق العمودي: جنين واحد وكرة من نسيج بطني.
\textbf{23.} وحدها: كرة بشروية. ومع الخلايا النباتية: أديم متوسط
--- عضل \omterm{prop:b2:vertebrate-development:gastrulation}{وحبل ظهري} --- مستحَث في القلنسوة.
\textbf{24.} لإثبات أن المحور الثاني استُحث في خلايا المضيف
ولم يبنه الطُّعم؛ فلو كان الطُّعم غير مميز،
لما أمكن استبعاد التمايز الذاتي. والطُّعم من
نوع مختلف ينجح (وقد استعملا نوعين من السمندل)، وهو ما
جعل الخلايا مميزة.
\textbf{25.} \omterm{def:b2:vertebrate-development:egg}{التفلج} الثاني عشر، و4096 خلية، عند نحو 6.75 ساعة؛ ونطق
من 38 و\qty{67}{\micro m}، وكل ما وراء \qty{105}{\micro
m}.
\end{solution}
